\section{Noise model, thresholds and measured false-detection levels} \label{app:decision}

\subsection{Noise model} \label{app:noise}

The noise model measures how deep a dip the light curve produces by chance at times where no transit is predicted.
From the cleaned light curve the predicted windows of both signals ($\pm 1.5\,\Tfourteen$) are cut. Fake phase offsets
of the tested signal are then placed between 0.15\,$P$ and 0.85\,$P$, in steps of $\Tfourteen/2$. For each fake phase
the light curve is detrended exactly as in Appendix~\ref{app:processing}, but with the fake ephemeris masked. The
windows at the fake times that pass the counting rule, with the segments and edge zones of that fake detrending, form
the noise pool. Every noise window therefore goes through the same detrending, masking, counting and measurement as a
real window. If the pool has fewer than 100 windows, the outcome is O4.

As a supplementary check, the whole set of counted windows at each fake phase gives a full-set statistic $d_\phi$, and
the number of fake phases with $d_\phi \geq d$ is reported. This number is not part of the outcome: it is in-sample, its
sets have their own number of windows, and, because the phases overlap, the count can only distinguish levels of about
5\,\%.

\subsection{Thresholds} \label{app:thresholds}

For each signal, with $n$ counted transits and statistic $d$:
\begin{itemize}
\item the noise pool is bootstrapped 20\,000 times to give the distribution of the mean of $n$ windows;
\item $\ddet$ is the 99.9th percentile of this distribution;
\item $\dlo$ is the 0.1st percentile of (noise $+$ a depth drawn from $N(D, D_{\rm err})$);
\item $\dhi$ is the 99.9th percentile of (noise $+$ a depth drawn from $N(D_{98}, D_{98,\rm err})$);
\item the sensitivity $\Sens = P(\mathrm{detect} \mid \mathrm{real})$ is the fraction of (noise $+$ $N(D, D_{\rm err})$)
  above $\ddet$.
\end{itemize}
All four quantities are computed for 20 bootstrap seeds, and their medians are used. The 99.9th percentile of 20\,000
draws rests on about 20 draws, so $\ddet$ moves between seeds (for B by up to about 10\,\%). The seed spread is
reported and plays no role in the outcome. A value of $d$ inside the spread is not described as ``marginal''.

\subsection{Measured false-detection levels} \label{app:fa}

The bootstrap of a finite pool of overlapping windows underestimates the tails of the noise, so the 99.9th percentile
is not the false-alarm level of the test. The levels that actually apply were therefore measured with a separate set of fake
ephemerides on the archival sectors: the true transits were removed, the fake periods were arbitrary, each fake phase
was one trial, and the thresholds were computed from the pool without the overlapping phases. The levels were measured
per product and per number of counted transits $n$ (Table~\ref{tab:fa}). The level that applies to Sector 110 is the
one of its product (SPOC 2-min for the primary run) and of the $n$ actually counted.

\begingroup\footnotesize
\begin{longtable}{>{\raggedright\arraybackslash}p{0.05\textwidth}>{\raggedright\arraybackslash}p{0.162\textwidth}>{\raggedright\arraybackslash}p{0.162\textwidth}>{\raggedright\arraybackslash}p{0.162\textwidth}>{\raggedright\arraybackslash}p{0.162\textwidth}>{\raggedright\arraybackslash}p{0.162\textwidth}}
\caption{Measured false-detection levels on archival data, from \code{calib/false\_alarm\_levels.json}: exceedances
$k$ out of $N$ trials. The first row pools all trials of a product and signal, with the 95\,\% interval from a bootstrap
over the trial ephemerides; the other rows give exactly $n$ counted transits. ``Bound'': one-sided 95\,\% Clopper-Pearson
upper bound for independent trials (see text). Empty: no trial with this $n$. The five columns contain all
17\,835 trials. The TESS-SPOC 30-min trials (Sector 6) are not quoted by \code{s110\_summary.py}: the Sector 110
TESS-SPOC product has 200-s cadence, and for it the SPOC 2-min levels are quoted.\label{tab:fa}}\\
\toprule
$n$ & SPOC 2-min (primary), A & SPOC 2-min (primary), B & QLP 200-s (fallback 1), A & QLP 200-s (fallback 1), B & TESS-SPOC 30-min (Sector 6), A \\
\midrule
\endfirsthead
\toprule
$n$ & SPOC 2-min (primary), A & SPOC 2-min (primary), B & QLP 200-s (fallback 1), A & QLP 200-s (fallback 1), B & TESS-SPOC 30-min (Sector 6), A \\
\midrule
\endhead
all & 0.47\,\% (26/5514; interval 0.27 to 0.70; bound 0.65\,\%) & 0.49\,\% (51/10312; interval 0.35 to 0.64; bound 0.62\,\%) & 0/614 (bound 0.49\,\%) & 0.64\,\% (5/777; interval 0.29 to 0.94; bound 1.35\,\%) & 0/618 (bound 0.48\,\%) \\
1 & 0/23 (bound 12.21\,\%) & 20/3054 (0.65\,\%, bound 0.95\,\%) & 0/1 (bound 95.00\,\%) & 2/188 (1.06\,\%, bound 3.31\,\%) & 0/11 (bound 23.84\,\%) \\
2 & 1/68 (1.47\,\%, bound 6.79\,\%) & 31/6040 (0.51\,\%, bound 0.69\,\%) & 0/6 (bound 39.30\,\%) & 2/325 (0.62\,\%, bound 1.92\,\%) & 0/51 (bound 5.70\,\%) \\
3 & 7/642 (1.09\,\%, bound 2.04\,\%) & 0/1077 (bound 0.28\,\%) & 0/33 (bound 8.68\,\%) & 1/202 (0.50\,\%, bound 2.33\,\%) & 0/140 (bound 2.12\,\%) \\
4 & 6/1538 (0.39\,\%, bound 0.77\,\%) & 0/141 (bound 2.10\,\%) & 0/136 (bound 2.18\,\%) & 0/62 (bound 4.72\,\%) & 0/175 (bound 1.70\,\%) \\
5 & 1/1203 (0.08\,\%, bound 0.39\,\%) &  & 0/197 (bound 1.51\,\%) &  & 0/118 (bound 2.51\,\%) \\
6 & 5/963 (0.52\,\%, bound 1.09\,\%) &  & 0/102 (bound 2.89\,\%) &  & 0/71 (bound 4.13\,\%) \\
7 & 6/907 (0.66\,\%, bound 1.30\,\%) &  & 0/66 (bound 4.44\,\%) &  & 0/51 (bound 5.70\,\%) \\
8 & 0/146 (bound 2.03\,\%) &  & 0/35 (bound 8.20\,\%) &  & 0/1 (bound 95.00\,\%) \\
9 & 0/24 (bound 11.73\,\%) &  & 0/32 (bound 8.94\,\%) &  &  \\
10 &  &  & 0/4 (bound 52.71\,\%) &  &  \\
11 &  &  & 0/2 (bound 77.64\,\%) &  &  \\
\bottomrule
\end{longtable}
\endgroup

Every trial set that was run is included: the three-product sets, two later SPOC 2-min sets, and the SPOC 2-min set run
by audit 3 (\code{build/false\_alarm\_check.py}, modes A, B, A\_S32, B\_S32, A\_S32\_a3 and B\_S32\_a3). Three of audit
3's ephemerides were recomputed by the trial script of the deposit and agree exactly. Since audit 6 every trial uses
exactly the thresholds of the Sector 110 run (\code{s110\_final.thresholds}, the median of 20 seeds). Before that, the
trial script used one seed, which changed 14 of the 17\,835 trial classifications (4 more and 10 fewer exceedances; all
trial depths identical).

``Bound'' is the one-sided 95\,\% Clopper-Pearson upper bound for independent trials \citep{ClopperPearson1934}. The
trials are not independent: neighbouring phases share data, and all trials of one product share one light curve. The
bound is therefore an approximation whose coverage is not established. It is not a guaranteed maximum, and neither the
bound nor the interval is a guarantee for Sector 110. These levels are frequencies of the noise model per trial, not
probabilities that a detected signal is a planet. At a given $n$, the level of B is not always higher than that of A. For an $n$ with few
trials the point value can be 0 while its bound is several times the overall level (A, $n = 8$: 0 of 146 trials, bound
2.03\,\%). The sentence of Appendix~\ref{app:wording} prints both, and the bound is the number to read. A grouping of $n$
that was tried after seeing the per-$n$ counts is not used. \code{s110\_summary.py} quotes, for the product of the run,
the overall level and the level for exactly the $n$ counted, each with its binomial bound and the remark on dependent
trials.

A check of the lower tail gave the matching numbers for real signals: a signal at the expected depth fell below $\dlo$
in 0.10\,\% (A) and 0.17\,\% (B) of the A\_S32 and B\_S32 trials. For fallback 2 (TESS-SPOC 200-s) no trials exist,
and the SPOC 2-min levels are quoted with that remark.

The trial files are \code{calib/false\_alarm\_*.json}, made by \code{build/false\_alarm\_check.py}; the levels are in
\code{calib/false\_alarm\_levels.json}, built by \code{build/false\_alarm\_levels.py}. The intervals come from a
bootstrap over the trial ephemerides. Because all trials of one product share one light curve, these intervals are
optimistic. The percentiles were set once (99.9 and 0.1) and are not tuned further. The levels are reported as they are,
and \code{s110\_summary.py} quotes the level of the product and of the $n$ of the actual run.

The transit times are fixed in advance and nothing is searched, so these are fixed-epoch levels, not search-wide
false-alarm rates. There is no minimum number of transits: the thresholds and the sensitivity are computed for the $n$
actually counted. An exclusion (O3) needs $\Sens \geq 0.95$ for that $n$. Whether this is reached depends on the
expected depth, the noise and $n$, not on a fixed number of transits; in the archival trials some sets reach it with
$n = 2$. A detection is possible for any $n$, and its false-detection level is the one of Table~\ref{tab:fa} for that
$n$.
